% !TeX root = ../../Book3.tex
\section{Carathéodory 边界对应}
\label{b3:sec:0702}

Jordan 区域的边界是一条简单闭曲线，此时每个边界点只对应一个素端。若边界仅局部连通，则可得到连续延拓，但边界映射仍可能多对一。三种陈述应当分别理解。

\subsection{Jordan 区域}
\label{b3:subsec:070201}

\term[jordan-quyu]{Jordan 区域}[Jordan domain]指以 Jordan 曲线为边界的有界内部区域。Jordan 曲线是一张圆周的连续单射像，不要求可微、有限长度或有切线。曲线的紧性使该参数化自动是同胚。多边形内部、圆盘及具有尖点但没有自接触的简单闭曲线内部都属于此类。

\ProseParagraphBreak
平面分离定理进一步给 Jordan 曲线两侧的拓扑圆盘结构。在 Jordan 区域中，以靠近一个边界点的两段边界弧和内部短弧构成横割，便能得到印象只含该点的链；反过来，每条零链的分支在边界圆周次序中挡住越来越短的弧。因此素端与普通边界点一一对应。这个平面拓扑事实不涉及全纯函数的导数，也不要求边界整齐。

\subsection{连续边界延拓}
\label{b3:subsec:070202}

\begin{theorem}[深层连续延拓判别]\label{b3:thm:07-continuous-extension}
设 \(U\) 有界且单连通，\(f:\mathbb D\rightarrow U\) 双全纯。则 \(f\) 有到 \(\overline{\mathbb D}\) 的连续延拓，当且仅当 \(\partial U\) 局部连通。
\end{theorem}
\begin{proofsketch}
本定理作为素端理论的进一步输入。局部连通性使每个素端印象为单点，并使印象投影连续，于是与素端对应复合得到延拓。反向利用连续圆周像的局部连通性；圆周参数化的连续性先给均匀小像，再按有限弧分割控制邻域中的连通片。一般边界上的精确拓扑论证见前引 Pommerenke 第2章。此处不把仅有印象单点的点态结论误当成全圆周连续性证明。
\end{proofsketch}
裂盘的边界局部连通，因此映射可连续延到圆周，但裂缝两侧对应不同的圆周点。连续延拓可以满射到整个边界，却不必单射。无界区域应在球面中作相应陈述；欧氏意义的有限边界值可能在无穷远点失效。

\subsection{边界同胚}
\label{b3:subsec:070203}

\term[caratheodory-bianjie-dingli]{Carathéodory 边界对应定理}[Carathéodory boundary correspondence theorem]在 Jordan 情形给出一一对应。

\begin{theorem}[Carathéodory 边界同胚]\label{b3:thm:07-caratheodory-jordan}
若 \(U\) 为 Jordan 区域，任何双全纯 \(f:\mathbb D\rightarrow U\) 唯一延为 \(\overline{\mathbb D}\rightarrow\overline U\) 的同胚。
\end{theorem}
\begin{proof}
Jordan 边界的素端投影是同胚，如本节第一小节的平面分离描述。将它与定理\ref{b3:thm:07-prime-end-correspondence}的素端延拓复合，得到闭圆盘到闭区域的连续双射。源空间紧、目标 Hausdorff，因此逆也连续。内部延拓的唯一性由内部点稠密得到。
\end{proof}
这是一条依赖已声明深层输入的完整推论证明，而不是对全部素端理论的替代。它不声称边界上的导数存在。例如角点附近的导数一般按一个幂次趋于0或无穷大；边界同胚仍然完全成立。

\begin{proposition}[圆边反射的局部形式]\label{b3:prop:07-circle-reflection}
若全纯函数 \(g\) 在圆的一侧邻域内无零，连续到一段圆弧，且该弧上 \(|g|=1\)，则可按反射公式跨弧延拓。对于单位圆，另一侧的值为
\[
 \widetilde g(z)=\frac1{\overline{g(1/\bar z)}}.
\]
\end{proposition}
\begin{proof}
反演 \(z\mapsto1/\bar z\) 交换单位圆两侧，复共轭和倒数把原侧全纯无零函数变成另一侧的全纯函数；圆弧上 \(|g|=1\) 使两个表达式相等。局部用一个 Möbius 变换把该圆弧拉直，再对目标单位圆作同样的拉直，转化为直径上取实值的反射。直边反射给拼接全纯，逆变换恢复所写公式。选择的小邻域避开 Möbius 变换的极点，且无零条件保证倒数合法。
\end{proof}
若圆弧是边界的一部分，这只延拓到该弧的局部另一侧，不保证映射仍在整个扩大区域单射。\(g(z)=z^m\) 的单位圆边界值模为1，确可全纯反射，但 \(m>1\) 时在盘内不是单叶映射。反射延拓与全局单射是两种不同结论。

\ProseParagraphBreak
光滑弧上的反射还给边界导数的意义。若源和目标边界在一点附近均为解析弧，且映射作为边界同胚把两段对应，拉直后得到两侧全纯映射。只要延拓后的零点阶数为1，就有非零边界导数；若目标有角点，则需先用幂次拉直，导数可趋于0或无穷。因此“边界延拓连续”不足以推出边界导数存在，更不足以推出导数不为零。

\subsection{素端表述}
\label{b3:subsec:070204}

素端表述在作边界标记时尤其有用。一个四边形需指定四个循环有序的边界素端，而不能仅给四个平面点；在裂缝区域，同一个点可有两个不同的标记。边界对应把这种标记传到圆盘圆周，随后可由 Möbius 变换把三个标记固定。剩下的第四个位置给出一个实际的共形参数。

\ProseParagraphBreak
若只关心内部映射，第5章不需要任何素端前提；若需要多边形边界值或矩形四角位置，就必须使用本章的边界同胚。这个依赖界面解释了为何 Schwarz--Christoffel 公式的证明虽然从一个内部 Riemann 映射开始，却还需要辨认实轴上的预顶点。

\ProseParagraphBreak
边界映射的三个层次可以在计算前作一次检查。第一层只有素端同胚，得到圆周参数，但一个参数点可能对应普通边界中的一整个印象。第二层有普通闭包上的连续映射，参数点有唯一边界值，但不同参数点仍可能对应同一点。第三层是 Jordan 边界同胚，参数点与普通边界点一一对应。Schwarz--Christoffel 的预顶点需要第三层，裂缝边界的两侧计算至少要保留第二层中的两份参数资料。若把裂缝的两岸在预顶点列表里合成一条边，就会丢掉循环次序。

\ProseParagraphBreak
Jordan 边界没有隐含弧长参数。某条简单闭曲线可以长度无穷，Carathéodory 定理仍给同胚，但沿它的分段光滑积分未必有定义。要写经典边界 Cauchy 积分，还需边界可求长或改用内部逼近围道。反射则需要解析边段，角点公式需要确切的局部角域；这些要求都比边界同胚更强。具体应用应当先决定所需的是拓扑值、路径积分还是导数，再列对应的边界条件。
